% sections/06_results.tex

\section{Results}
\label{sec:results}

All numeric cells in this section are placeholders populated from the
Layer-1 harness.
Run the harness for each scenario (e.g.\ \texttt{python scripts/run\_evaluation.py --scenario s1 --seed 42}; repeat for s1..s4)
and replace each \texttt{\textbackslash TODO\{...\}} with the printed value.

\subsection{Naive vs.\ Grouped Cross-Validation Gap}
\label{sec:cv_gap}

Table~\ref{tab:cv_gap} reports the mean accuracy (classification) or $R^2$
(regression) under naive stratified KFold versus group-aware
GroupKFold / LeaveOneGroupOut, for each scenario.
Figure~\ref{fig:cv_gap} visualizes the per-scenario gap as a bar chart.
A positive gap indicates that naive CV overestimates performance due to
replicate leakage.

\begin{table}[ht]
  \centering
  \caption{Naive KFold vs.\ grouped CV performance gap. Higher gap = more leakage.}
  \label{tab:cv_gap}
  \begin{tabular}{lrrrr}
    \toprule
    Scenario & Metric & Naive CV & Grouped CV & Gap \\
    \midrule
    S1 (wear-layer) & $R^2$    & TODO & TODO & TODO \\
    S2 (species)    & Accuracy & TODO & TODO & TODO \\
    S3 (vinyl/wood) & Accuracy & TODO & TODO & TODO \\
    S4 (FTIR purity)& Accuracy & TODO & TODO & TODO \\
    \midrule
    Mean gap        &          &      &      & \TODO{mean gap across scenarios} \\
    \bottomrule
  \end{tabular}
\end{table}

\begin{figure}[ht]
  \centering
  \IfFileExists{figures/fig2_cv_gap.pdf}{%
    \includegraphics[width=0.75\linewidth]{figures/fig2_cv_gap}%
  }{%
    \fbox{\parbox{0.9\linewidth}{\centering\vspace{2cm}\TODO{figure: fig2\_cv\_gap}\vspace{2cm}}}%
  }
  \caption{
    Naive KFold vs.\ grouped CV accuracy/R\textsuperscript{2} by scenario.
    Error bars show inter-fold standard deviation.
    \textbf{[TODO: generate from harness run]}
  }
  \label{fig:cv_gap}
\end{figure}

\subsection{Leakage Detection Scorecard}
\label{sec:leakage_results}

Table~\ref{tab:leakage} reports TPR, FPR, and precision of the leakage
detector on the 6 test fixtures (3 positive, 3 clean).

\begin{table}[ht]
  \centering
  \caption{Leakage detection performance on 6 fixtures (3 positive, 3 clean).
    Positive = leakage present; clean = no leakage.}
  \label{tab:leakage}
  \begin{tabular}{lrrr}
    \toprule
    Check type & TPR & FPR & Precision \\
    \midrule
    replicate\_leakage    & \TODO{TPR} & \TODO{FPR} & \TODO{Precision} \\
    group\_leakage\_risk  & \TODO{TPR} & \TODO{FPR} & \TODO{Precision} \\
    Combined              & \TODO{TPR} & \TODO{FPR} & \TODO{Precision} \\
    \bottomrule
  \end{tabular}
\end{table}

\subsection{HITL Compliance, Fallback, and Plan Quality}
\label{sec:hitl_results}

Table~\ref{tab:hitl} summarizes HITL compliance, fallback correctness,
and plan quality auto-score across the four scenarios.

\begin{table}[ht]
  \centering
  \caption{HITL compliance, fallback correctness, and plan quality auto-score.}
  \label{tab:hitl}
  \begin{tabular}{lrrrr}
    \toprule
    Scenario & HITL Compliance & Fallback Correct & Plan Auto-score & Plan Auto-pass \\
    \midrule
    S1 & TODO & TODO & TODO & TODO \\
    S2 & TODO & TODO & TODO & TODO \\
    S3 & TODO & TODO & TODO & TODO \\
    S4 & TODO & TODO & TODO & TODO \\
    \midrule
    Mean & \TODO{mean} & \TODO{mean} & \TODO{mean} & \TODO{mean} \\
    \bottomrule
  \end{tabular}
\end{table}

\subsection{Ablation Deltas}
\label{sec:ablation}

Table~\ref{tab:ablation} reports the delta in leakage TPR, HITL compliance,
and plan quality auto-score for each ablation configuration relative to the
full system.
Negative deltas indicate degradation.

\begin{table}[ht]
  \centering
  \caption{Ablation deltas relative to the full system.
    $\Delta$TPR: change in leakage detection TPR.
    $\Delta$HITL: change in HITL compliance rate.
    $\Delta$Plan: change in plan quality auto-score.}
  \label{tab:ablation}
  \begin{tabular}{lrrr}
    \toprule
    Configuration       & $\Delta$TPR & $\Delta$HITL & $\Delta$Plan \\
    \midrule
    Full (baseline)     & 0 & 0 & 0 \\
    no\_guardrails      & TODO & TODO & TODO \\
    no\_validation      & TODO & TODO & TODO \\
    no\_memory          & TODO & TODO & TODO \\
    deterministic\_off  & TODO & TODO & TODO \\
    \bottomrule
  \end{tabular}
\end{table}

\subsection{Effort}
\label{sec:effort}

Each full-pipeline scenario run (inspect $\to$ report) requires 7 tool calls
(one per pipeline stage).
Wall-clock time per scenario at seed = 42 on a laptop CPU is
\TODO{per-scenario wall-clock seconds}.
No HITL decision points are incurred in Layer-1 (the harness passes approval
automatically); Layer-2 HITL effort is future work.

\subsection{Case Study: FTIR Purity (Antigravity-Final Run)}
\label{sec:casestudy}

To illustrate the pipeline behavior on a real FTIR dataset, we describe the
\texttt{antigravity-ftir-purity-run-final} run, which was executed outside the
Layer-1 harness using the MCP server with a human operator approving each gate.

The run executed 30 model--preprocessing combinations: 6 preprocessing methods
(SNV, MSC, first-derivative Savitzky--Golay, second-derivative Savitzky--Golay,
baseline correction, area normalization) $\times$ 5 model families
(SVM-RBF, Random Forest, Logistic Regression, XGBoost, and PCA-LDA).
The PCA-LDA models failed across all 6 preprocessing conditions due to
singular covariance matrices on this small-$n$ dataset; these were recorded in
\texttt{failed\_models} and surfaced as fallback candidates.

The best-measured accuracy was \textbf{0.905} (SVM-RBF or Random Forest with
the best preprocessing; exact model details available in the run artifact at
\texttt{runs/antigravity-ftir-purity-run-final/artifacts/run\_summary.json}).
Despite this strong accuracy, the validation step set
\texttt{validation\_passed = False} because 9 validation warnings were raised,
including small-sample-per-class warnings (the $\approx$20-sample, 7-group
dataset makes per-class counts sparse) and class-imbalance flags.

Consequently, \texttt{select\_best\_model} set
\texttt{approval\_required = True}, correctly routing the selection to the
human operator rather than autonomously committing to the 0.905 model.
This case study demonstrates the core HITL design: a high metric does not
suppress validation warnings, and validation warnings do not suppress the
metric---both are surfaced, and the human decides.
